% from ega3-6-fr.tex
\oldpage[III]{169}
On a alors, en posant
$\mathscr P_\bullet'^{(i)}=
\mathscr P_\bullet^{(i)}\otimes_{\mathscr O_S}\mathscr O_{S'}$, des
isomorphismes canoniques fonctoriels
\[
\begin{split}
  &\mathfrak{Tor}_n^{S'}
  (f'_1,\ldots,f'_m,1_{S'};
   \mathscr P_\bullet'^{(1)},\ldots,
   \mathscr P_\bullet'^{(m)},\mathscr Q'_\bullet)
  \xrightarrow{\sim}\\
  &\qquad
  \bigoplus_{n'+n''=n}
  \mathfrak{Tor}_{n'}^S
  (f_1,\ldots,f_m;
   \mathscr P_\bullet^{(1)},\ldots,
   \mathscr P_\bullet^{(m)})
  \otimes_{\mathscr O_S}\mathscr H_{n''}(\mathscr Q'_\bullet).
\end{split}
\tag{6.9.6.1}\label{III.6.9.6.1.fr}
\]
En particulier, pour $\mathscr Q'_\bullet$ réduit à un seul terme
$\mathscr F'$ de degré $0$, on a des isomorphismes canoniques fonctoriels
\[
\begin{split}
  &\mathfrak{Tor}_n^{S'}
  (f'_1,\ldots,f'_m,1_{S'};
   \mathscr P_\bullet'^{(1)},\ldots,
   \mathscr P_\bullet'^{(m)},\mathscr F')
  \xrightarrow{\sim}\\
  &\qquad
  \mathfrak{Tor}_n^S
  (f_1,\ldots,f_m;
   \mathscr P_\bullet^{(1)},\ldots,
   \mathscr P_\bullet^{(m)})
  \otimes_{\mathscr O_S}\mathscr F'
\end{split}
\tag{6.9.6.2}\label{III.6.9.6.2.fr}
\]
et plus particulièrement, pour $\mathscr F'=\mathscr O_{S'}$,
\[
\begin{split}
  &\mathfrak{Tor}_n^{S'}
  (f'_1,\ldots,f'_m;
   \mathscr P_\bullet'^{(1)},\ldots,
   \mathscr P_\bullet'^{(m)})
  \xrightarrow{\sim}\\
  &\qquad
  \mathfrak{Tor}_n^S
  (f_1,\ldots,f_m;
   \mathscr P_\bullet^{(1)},\ldots,
   \mathscr P_\bullet^{(m)})
  \otimes_{\mathscr O_S}\mathscr O_{S'}.
\end{split}
\tag{6.9.6.3}\label{III.6.9.6.3.fr}
\]
\end{corollary}

\begin{proof}
L'hypothèse de platitude sur les Modules composant les
$\mathscr P_\bullet^{(i)}$ entraîne que les complexes
$\mathscr{T}\!or_{q_i}^S(\mathscr P_\bullet^{(i)},\mathscr O_{S'})$ sont nuls
pour $q_i\ne0$ (6.5.8). La suite $(f')$ est donc dégénérée; l'hypothèse
$2^\circ$ entraîne d'ailleurs qu'elle est birégulière (6.9.3), donc le
edge-homomorphisme
\[
\begin{split}
  &\mathfrak{Tor}_n^S
  (f_1,\ldots,f_m,1_{S'};
   \mathscr P_\bullet^{(1)},\ldots,
   \mathscr P_\bullet^{(m)},\mathscr Q'_\bullet)
  \longrightarrow {}^{(f')}\mathscr E_{n0}^2\\
  &\qquad=
  \mathfrak{Tor}_n^{S'}
  (f'_1,\ldots,f'_m,1_{S'};
   \mathscr P_\bullet'^{(1)},\ldots,
   \mathscr P_\bullet'^{(m)},\mathscr Q'_\bullet)
\end{split}
\tag{6.9.6.4}\label{III.6.9.6.4.fr}
\]
est bijectif (0, 11.1.6). L'hypothèse de platitude sur les Modules
\[
  \mathfrak{Tor}_n^S
  (f_1,\ldots,f_m;
   \mathscr P_\bullet^{(1)},\ldots,\mathscr P_\bullet^{(m)})
\]
entraîne que ${}^{(e)}\mathscr E_{pq}^2=0$ pour $p\ne0$ (6.5.8). La suite
$(e)$ est donc aussi dégénérée, et comme elle est birégulière, le
edge-homomorphisme
\[
\begin{split}
  &\mathfrak{Tor}_n^S
  (f_1,\ldots,f_m,1_{S'};
   \mathscr P_\bullet^{(1)},\ldots,
   \mathscr P_\bullet^{(m)},\mathscr Q'_\bullet)
  \longrightarrow {}^{(e)}\mathscr E_{0n}^2\\
  &\qquad=
  \bigoplus_{n'+n''=n}
  \mathfrak{Tor}_{n'}^S
  (f_1,\ldots,f_m;
   \mathscr P_\bullet^{(1)},\ldots,
   \mathscr P_\bullet^{(m)})
  \otimes_{\mathscr O_S}\mathscr H_{n''}(\mathscr Q'_\bullet)
\end{split}
\tag{6.9.6.5}\label{III.6.9.6.5.fr}
\]
est bijectif (0, 11.1.6); d'où, en combinant les deux isomorphismes
précédents, l'isomorphisme (6.9.6.1). L'isomorphisme (6.9.6.2) s'en déduit
trivialement, puisque l'on a alors
$\mathscr H_q(\mathscr Q'_\bullet)=0$ si $q\ne0$ et
$\mathscr H_0(\mathscr Q'_\bullet)=\mathscr F'$. Enfin, le cas
$\mathscr F'=\mathscr O_{S'}$ dans (6.9.6.2) donne l'isomorphisme (6.9.6.3),
compte tenu de (6.7.12).
\end{proof}

\begin{corollary}[6.9.7]
Sous les conditions de (6.9.1), supposons $S$ et $S'$ affines, et supposons
donné pour chaque $i$ un entier $d_i$ ($1\leq i\leq m$). Il existe alors un
entier $N$ ne dépendant que de $S$, des $X^{(i)}$ et des $d_i$, ayant la
propriété suivante : pour tout entier $n_0$, on a des isomorphismes canoniques
(6.9.6.3) pour $n\leq n_0$ et pour tout système de complexes
$\mathscr P_\bullet^{(i)}$ vérifiant les conditions suivantes :
\begin{enumerate}
  \item[$1^\circ$] $\mathscr P_k^{(i)}=0$ pour $k<d_i$;
  \item[$2^\circ$] $\mathscr P_k^{(i)}$ est plat sur $S$ pour $k<n_0+N$;
  \item[$3^\circ$]
  $\mathfrak{Tor}_q^S(f_1,\ldots,f_m;
  \mathscr P_\bullet^{(1)},\ldots,\mathscr P_\bullet^{(m)})$ est plat sur
  $S$ pour $q<n_0+N$.
\end{enumerate}
\end{corollary}

\begin{proof}
Supposons $\mathscr P_k^{(i)}$ plat sur $S$ pour $k<r$; alors
$\mathscr{T}\!or_{q_i}^S(\mathscr P_k^{(i)},\mathscr O_{S'})=0$ pour $k<r$
et $q_i\ne0$; calculons
\[
\begin{split}
  \mathfrak{Tor}_p^{S'}\Bigl(
  &f'_1,\ldots,f'_m;
  \mathscr{T}\!or_{q_1}^S(\mathscr P_\bullet^{(1)},\mathscr O_{S'}),
  \ldots,\\
  &\mathscr{T}\!or_{q_m}^S(\mathscr P_\bullet^{(m)},\mathscr O_{S'})
  \Bigr).
\end{split}
\tag{6.9.7.1}\label{III.6.9.7.1.fr}
\]

